% !TeX program = xelatex
% UTF-8 transcription of all seven handwritten source pages, in page order.
% Original wording, abbreviations, values and arrows are retained.
% images/pic01.jpg = 20261012_184501556.jpg
% images/pic02.jpg = 20261012_184509620.jpg
% images/pic03.jpg = 20261012_184515761.jpg
% images/pic04.jpg = 20261012_184521990.jpg
% images/pic05.jpg = 20261012_184527894.jpg
% images/pic06.jpg = 20261012_184535892.jpg
% images/pic07.jpg = 20261012_184546555.jpg
% Standalone document; also usable with \input in an Xepersian document.
\ifdefined\settextfont
  \def\CHDStandalone{0}
\else
  \def\CHDStandalone{1}
  \documentclass[12pt,a4paper]{article}
  \usepackage[margin=2cm]{geometry}
  \usepackage{amsmath,amssymb,enumitem}
  \usepackage{xepersian}
  \IfFontExistsTF{Amiri}{\settextfont{Amiri}}{\settextfont{FreeSerif}}
  \setlatintextfont{FreeSerif}
\fi
\providecommand{\en}[1]{\lr{#1}}
\newcommand{\chdpage}[3]{\clearpage\section*{صفحهٔ #1}
  \noindent موضوع: #2\par
  {\footnotesize\en{#3}\par}}
\ifnum\CHDStandalone=1 \begin{document}\fi

\chdpage{۱}{براونوالد ـ ۸۲ ـ بیماری‌های مادرزادی قلبی در بزرگسالان}{pic01.jpg}
\subsection*{شانت‌ها چپ به راست}
\begin{itemize}
\item شانت‌های پروگزیمال به تری‌کاسپید (مثل \en{ASD}) موجب دیلاتاسیون قلب راست می‌شوند.
\item شانت‌های دیستال به تری‌کاسپید (\en{VSD, PDA}) موجب دیلاتاسیون قلب چپ می‌شوند.
\item با افزایش \en{PVR} (مقاومت عروق پولمونری)، شانت‌های بزرگ می‌توانند موجب تغییر جهت جریان خون از فرم چپی (راست به چپ) به \en{$L\to R$} شوند.
% The source explicitly writes «از فرم چپی (راست به چپ) به L→R».
\item سیانوز و سندرم آیزن‌منگر ایجاد می‌کند.
\end{itemize}
\subsection*{ارزیابی شانت‌ها}
\begin{itemize}
\item تعیین مقاومت عروق پولمونری (\en{PVR}) و سیستمیک (\en{SVR}).
\item تعیین نسبت جریان خون پولمونری (\en{$Q_p$}) به سیستمیک (\en{$Q_s$}).
\end{itemize}
\subsection*{\en{Atrial Septal Defects (ASDs) and Partial Anomalous Pulmonary Veins (PAPVC)}}
\begin{itemize}
\item یافته‌های بالینی: عدم تحمل فعالیت، تپش قلب، \en{DoE}.
\item آریتمی‌های فوق بطنی در ۱۰\% افراد بالای ۴۰ سال این گروه رخ می‌دهد.
\end{itemize}

\chdpage{۲}{براونوالد ـ ۸۲ ـ بیماری‌های مادرزادی}{pic02.jpg}
\begin{itemize}
\item وجود سیانوز مطرح‌کنندهٔ تغییر جهت شانت (\en{shunt reversal}) و سندرم آیزن‌منگر است.
\item برای ارزیابی شدت شانت: پالس اکسیمتری حین استراحت و فعالیت.
\item \en{Wide, fixed Splitting of the second heart sound}.
\item به طور کلی [خوانش نامطمئن: به عمده] اختلال قلبی مادرزادی زیر وجود داشت، فعالیت فیزیکی شدید باید اجتناب شود:
% The unclear connecting phrase is flagged rather than silently reconstructed.
  \begin{itemize}
  \item سیانوز.
  \item آنومالی‌های کرونری \en{high Risk}.
  \item هایپرتروفی.
  \item دیلاتاسیون شدید آئورت (آئورت به صورت حاشیه‌نویسی آمده است).
  \item انسداد شدید \en{outflow tract}.
  \item \en{PH} شدید.
  \item اختلال عملکرد بطنی شدید.
  \item آریتمی بطنی.
  \end{itemize}
\end{itemize}
\subsection*{\en{chamber enlargement}}
\begin{itemize}
\item \en{ASD: RA (+), RV (+), PA (+)}؛ \en{PA}: شریان ریوی.
\item \en{VSD: RV ($\pm$), PA (+), LA (+), LV (+)}.
\item \en{PDA: PA (+), LA (+), LV (+), Aorta (+)}.
\end{itemize}
\begin{itemize}
\item در \en{EKG}: \en{Incomplete RBBB, RAD}.
\item در \en{CXR}: کاردیومگالی، دیلاتاسیون شریان‌های پولمونری سنترال، کانجسشن ریوی.
\end{itemize}

\chdpage{۳}{براونوالد ـ ۸۲ ـ بیماری‌های مادرزادی}{pic03.jpg}
\begin{itemize}
\item تصویربرداری از قلب ضروری است.
\item اولین تصویربرداری \en{TTE} است.
\item ممکن است \en{TEE, CMR, CCT}.
\end{itemize}
\subsection*{کدام \en{ASD} را ترمیم می‌کنیم؟}
\begin{itemize}
\item هم جراحی، هم ترنس \en{Cath}.
\item وجود شواهد دایلیشن قلب راست که همراه با علامت باشد.
\end{itemize}
\subsection*{کرایتریا}
{\footnotesize در عنوانِ تصویر، پس از «کرایتریا»: [واژهٔ خط‌خورده، ناخوانا].\par}
\begin{enumerate}
\item \en{$\mathrm{PVR}<\frac{1}{3}\mathrm{SVR}$}.
\item \en{$\mathrm{PA\ BP}<\frac{1}{2}\,\text{systemic BP}$}.
\item انلارجمنت قلب راست.
\item \en{$Q_p/Q_s$} حداقل \en{$1.5:1$} باشد.
\end{enumerate}
\begin{itemize}
\item اگر صرفاً شواهد انلارجمنت قلب راست داشته باشیم ولی علامت بالینی نداشته باشیم باز هم توصیه به ترمیم می‌شود. \en{COR IIa}.
\item بهترین زمان ترمیم \en{ASD} قبل از ۲۵ سالگی است.
\item ترمیم \en{ASD} تأثیری در کاهش خطر بروز آریتمی ندارد.
\end{itemize}

\chdpage{۴}{براونوالد ـ ۸۲ ـ بیماری‌های مادرزادی}{pic04.jpg}
\subsection*{پس از ترمیم کدام بیماران باید با \en{TTE} به صورت دوره‌ای بررسی شوند؟}
\begin{itemize}
\item باقی ماندن مقدار قابل توجهی از شانت.
\item اختلال عملکرد دریچه‌ای، اختلال عملکرد بطنی.
\item آریتمی.
\item \en{PH}.
\end{itemize}
\begin{itemize}
\item پس از همهٔ ترمیم‌ها: \en{TTE}، ۱ هفته، ۱ ماه، سپس سالانه برای حداقل ۵ سال.
\end{itemize}
{\footnotesize یادداشت خط‌خورده و دایره‌کشیده در تصویر: \en{VSD}.\par}
\subsection*{\en{Atrioventricular Septal Defects (AVSD)}}
\begin{itemize}
\item طیفی از ضایعات هستند که شامل نارسایی‌های سپتوم \en{AV} که معمولاً با اختلالات دریچه‌های \en{AV} همراه هستند می‌شوند.
\item \en{Complete AVSD}: هر دو سپتوم دهلیزی و بطنی دیفکت وجود دارند و دریچهٔ \en{AV} معمولاً ۵ لتی است.
\item \en{Partial AVSD}: جزء \en{VSD} را ندارد و تقریباً همیشه در روی لت قدامی دریچهٔ میترال یک \en{cleft} وجود دارد.
\item علائم بالینی: معمولاً در کودکی ترمیم می‌شود و به بزرگسالی نمی‌رسد.
\end{itemize}

\chdpage{۵}{براونوالد ـ ۸۲ ـ بیماری‌های مادرزادی}{pic05.jpg}
\begin{itemize}
\item اگر \en{AVSD} تا بزرگسالی ترمیم نشده باشد:
  \begin{itemize}
  \item معمولاً بیماری برگشت‌ناپذیر عروق پولمونری پیدا می‌کند.
  \item سندرم آیزن‌منگر شایع است.
  \end{itemize}
\item فرم شایع‌تر دیگری در بزرگسالان به این صورت است که در کودکی ترمیم شده ولی الان با شواهد رگورژیتاسیون دریچهٔ \en{AV}، استنوز دریچه، انسداد \en{LVOT} یا آریتمی دهلیزی مراجعه کرده است.
\item اندیکاسیون‌های ترمیم \en{AVSD} در بزرگسالان مشابه \en{ASD} است.
\item صرفاً جراحی؛ ترنس \en{Cath} نداریم.
\item درصد قابل توجهی دچار بیماری دریچهٔ میترال و انسداد \en{LVOT} به دنبال ترمیم در کودکی می‌شوند.
\item \en{Complete HB} یک عارضهٔ دیگر آن است: پس از ترمیم.
\item بیمار با ترمیم \en{AVSD} با هر درجه‌ای از \en{PH} باید به صورت دوره‌ای پیگیری شود.
\end{itemize}

\chdpage{۶}{براونوالد ـ ۸۲ ـ بیماری‌های مادرزادی (\en{CHD})}{pic06.jpg}
\subsection*{\en{VSDs}}
\begin{itemize}
\item شایع‌ترین نوع \en{CHD} است.
\item علائم بالینی:
  \begin{itemize}
  \item سوفل هولوسیستولیک هارش.
  \item سندرم آیزن‌منگر.
  \end{itemize}
\end{itemize}
\subsection*{اندیکاسیون ترمیم}
\begin{itemize}
\item شواهد والیوم اورلود بطن چپ + شانت از نظر همودینامیک \en{Signif} باشد: \en{$Q_p/Q_s>1.5/1$}.
\item استاندارد طلایی جراحی مستقیم ولی ترنس \en{Cath} هم اوکی.
\item برای ترمیم باید:
  \begin{itemize}
  \item \en{$\mathrm{PA\ SBP}<\frac{1}{2}\,\text{systemic SBP}$}.
  \item \en{$\mathrm{PVR}<\frac{1}{3}\mathrm{SVR}$}.
  \end{itemize}
\end{itemize}
\subsection*{اندیکاسیون پیگیری طولانی‌مدت بیمار پس از ترمیم \en{VSD}}
\begin{itemize}
\item هر نوع اختلال دریچه‌ای.
\item اختلال عملکرد بطنی.
\item آریتمی.
\item \en{PH}.
\end{itemize}
\subsection*{\en{PDA}}
\begin{itemize}
\item ارتباط \en{PA} به آئورت نزولی؛ ترمیم ترنس \en{Cath}.
\item علائم بالینی: \en{CHF, PH}، سندرم آیزن‌منگر.
\end{itemize}

\chdpage{۷}{براونوالد ـ ۸۲ ـ \en{CHD} در بزرگسالان}{pic07.jpg}
\begin{itemize}
\item \en{$\mathrm{SpO}_2$} در هر چهار اندام باید بررسی گردد.
\item برای بررسی شانت راست به چپ (\en{$R\to L$}) و سیانوز.
\item بزرگسالان با \en{PDA} جدی معمولاً \en{Wide PP} دارند.
\item در سمع سوفل در زیر کلاویکل چپ سمع می‌شود.
\end{itemize}
\subsection*{اندیکاسیون ترمیم \en{PDA} در بزرگسالان}
\begin{itemize}
\item \en{LAE} یا \en{LVE} وجود داشته باشد که آن را مرتبط با \en{PDA} بدانیم.
\item شانت خالص \en{$L\to R$}.
\item \en{$\mathrm{PA\ SBP}<\frac{1}{2}\,\text{Systemic SBP}$}.
\item \en{$\mathrm{PVR}<\frac{1}{3}\mathrm{SVR}$}.
\end{itemize}
\begin{itemize}
\item بیمارانی که پس از ترمیم دچار انسداد آئورت یا شریان ریوی چپ (\en{LPA}) شده‌اند باید به طور دوره‌ای تصویربرداری انجام دهند.
\item پروفیلاکسی اندوکاردیت برای ۶ ماه در ترمیم، و تا آخر عمر اگر \en{PDA} باقی ماند.
\item \en{PDA} که نمی‌خواهیم ترمیم کنیم نیاز به پروفیلاکسی ندارد.
\end{itemize}
% Printed form labels: موضوع on all seven pages; تاریخ on pages 2--7,
% with blank date fields. Page 1 handwriting overlaps the تاریخ field.
% Original handwritten page numbers: ۱، ۲، ۳، ۴، ۵، ۶، ۷.
% Each page has a small printed notebook logo without readable medical text.
\ifnum\CHDStandalone=1
  \let\CHDEnd\enddocument
\else
  \let\CHDEnd\relax
\fi
\CHDEnd
